%% CRML v5 — IEEE TNNLS Submission
%% Cross-Relational Multi-Scale Learning with Hypergraph Neural Networks
%% for Stock Return Prediction
\documentclass[journal]{IEEEtran}

\usepackage{amsmath}
\usepackage{amsfonts}
\usepackage{amssymb}
\usepackage{graphicx}
\usepackage{cite}
\usepackage{algorithm}
\usepackage{algpseudocode}
\usepackage{placeins}
\usepackage{booktabs}
\usepackage{tikz}
\usetikzlibrary{positioning, arrows.meta, shapes.geometric, fit, backgrounds}

\renewcommand{\topfraction}{0.85}
\renewcommand{\bottomfraction}{0.85}
\renewcommand{\textfraction}{0.10}
\renewcommand{\floatpagefraction}{0.85}
\setcounter{topnumber}{4}
\setcounter{bottomnumber}{4}
\setcounter{totalnumber}{6}

\makeatletter
\newcommand{\DescState}[1]{%
  \State \parbox[t]{\dimexpr\linewidth-\algorithmicindent}{#1\strut}}
\makeatother

\begin{document}

\title{Cross-Relational Multi-Scale Learning with Hypergraph Neural Networks
       for Stock Return Prediction}

\author{\IEEEauthorblockN{Anmol Jain, Aditya Ranjan Patra, and Sibarama Panigrahi}
\IEEEauthorblockA{Department of Computer Science and Engineering\\
National Institute of Technology Rourkela, Odisha, India}}

\IEEEtitleabstractindextext{%
\begin{abstract}
Cross-sectional stock return prediction is a ranking problem: the
question is not whether the market rises, but which stocks outperform
their peers. Existing graph-based methods encode inter-stock
dependencies at a single structural scale, either through pairwise
edges or fixed group memberships. This paper asks whether jointly
modeling all three dominant relational scales in equity markets---static
sector co-membership, dynamic similarity clustering, and directed
pairwise lead-lag---can generate statistically significant cross-sectional
signal under a minimal feature set; whether each scale contributes
independently; and whether the resulting architecture transfers across
equity markets with different microstructural characteristics. We
propose CRML (Cross-Relational Multi-Scale Learning), which constructs
three structures in parallel: a static GICS sector hypergraph, a dynamic
dual-polarity KNN hypergraph connecting each stock to its $k$ most similar
co-movers and $k$ least similar inverse movers through separate hyperedge
channels, and a GATv2 pairwise attention graph. A gating layer fuses
these with per-stock softmax weights recomputed on each forward pass.
On the S\&P~100 with 12 technical features and a 5-day simple return
target $r{=}(C_{t+5}/C_t){-}1$, evaluated over four held-out walk-forward
test folds with hyperparameter search confined to earlier folds, CRML
achieves mean daily IC~$=0.047$ (95\,\% bootstrap CI: $[0.016,\,0.077]$),
against 0.020 for ridge regression. Ablation confirms independent
contributions from each scale; sector hypergraph removal causes the
largest degradation ($\Delta\mathrm{IC}{=}{-}0.012$, $p{=}0.053$).
Replacing the training objective with the standard Qlib MSE-only
protocol collapses IC by 84\,\%, showing that benchmark training recipes
do not transfer to dual-head architectures. On the NIFTY~100, no
statistically significant signal is found (IC~$=0.016$,
CI $[-0.015,\,0.039]$), identifying structural conditions under which
this class of architecture fails.
\end{abstract}

\begin{IEEEkeywords}
Graph Neural Networks, Hypergraph Learning, Stock Return Prediction,
Quantitative Finance, Relational Learning, Deep Learning.
\end{IEEEkeywords}}

\maketitle
\IEEEdisplaynontitleabstractindextext
\IEEEpeerreviewmaketitle

\section{Introduction}

\IEEEPARstart{E}{quity} markets price hundreds of stocks simultaneously,
and the return of any single stock on a given day is not determined in
isolation. A regulatory change affecting semiconductor supply chains can
move Apple, NVIDIA, and Broadcom within hours of each other---not because
of any direct financial link between them, but because they share
exposure to the same production bottleneck. A risk-off episode triggered
by an unexpected inflation print compresses high-beta growth stocks
across every sector simultaneously, while defensive utilities and
consumer staples decouple from the broader selloff. When Microsoft
reports earnings, smaller software peers re-rate in the same session as
investors update their sector-wide expectations. These patterns are not
noise. They reflect structural, economically motivated inter-stock
dependencies that carry cross-sectional predictive
signal~\cite{jegadeesh1993returns}---and any model treating stocks as
independent time series discards that signal entirely.

The prediction problem we study is cross-sectional: given the feature
histories of $N$ stocks on day $t$, the goal is to rank the universe
by expected forward return over the next five trading days. This differs
from univariate forecasting in a concrete way. A model that predicts
every stock will decline 2\,\% during a broad market selloff is a
perfect time-series forecaster if the index drops exactly 2\,\%---yet
it has zero cross-sectional utility. It cannot separate the stock that
falls 5\,\% from the one that falls 0.5\,\%. Our evaluation metric is
the Information Coefficient~(IC)~\cite{grinold2000active}, the Spearman
rank correlation~\cite{spearman1904} between predicted and realized
returns computed per trading day. Positive IC supports a long-short
portfolio insulated from market beta~\cite{grinold2000active}.

Graph neural networks have become the standard approach for encoding
inter-stock dependencies~\cite{patel2024gnn}. Every recent architecture,
however, commits to a single topological representation.
RSR~\cite{feng2019temporal} and HATS~\cite{kim2019hats} use pairwise
edges from static knowledge graphs. STHAN-SR~\cite{sawhney2021sthan}
replaces these with a single sector-defined hypergraph.
MASTER~\cite{li2024master} deploys full transformer attention over all
stock-time pairs. HIST~\cite{xu2021hist}, the most expressive
single-scale predecessor, combines a predefined industry concept graph
with a learned hidden concept graph, but both modules remain at the
group level and neither models directed pairwise effects. These
architectures fix their relational mechanism at design time and apply
it uniformly across all stocks and regimes.

That constraint matters because the three dominant relational
mechanisms in equity markets operate at different structural scales.
Sector co-membership is stable: GICS classifications reflect durable
economic realities~\cite{msci2023gics}, and sector co-movement persists
across regimes~\cite{jegadeesh1993returns,gu2020empirical}. Dynamic
similarity is transient: during a macro shock, unrelated stocks can
lock into cross-sector correlations for days before decoupling.
Pairwise lead-lag is directed: a megacap bellwether asymmetrically
predicts a smaller peer's return, not the reverse. A static sector
hypergraph cannot capture transient cross-sector clustering; a dynamic
KNN graph discards directional asymmetry; a pairwise attention graph
loses the unified group signal when an entire sector rotates. None
subsumes the others.

We introduce CRML, which constructs all three relational structures in
parallel and combines them through learned, input-conditioned gates.
Each stock's 60-day feature history is compressed by a shared GRU
encoder~\cite{cho2014learning}. The embeddings feed simultaneously into
a static GICS sector hypergraph via HGNN
convolution~\cite{feng2019hypergraph}, a dynamic dual-polarity KNN
hypergraph connecting each stock to its $k$ most similar co-movers and
$k$ least similar inverse movers through separate hyperedge channels
rebuilt on every forward pass, and a sparse GATv2 pairwise attention
graph~\cite{brody2022attentive}. An input-conditioned gating layer
computes per-stock, per-day softmax weights over the three outputs. The
architecture is evaluated on the S\&P~100~(US) and NIFTY~100~(India)
under a strict walk-forward protocol~\cite{lopezdeprado2018advances}
that separates hyperparameter optimization from testing.

This paper addresses three research questions.
\textbf{RQ1:} Can jointly modeling sector co-membership, dynamic
similarity clustering, and directed pairwise attention generate a
statistically significant cross-sectional signal on a highly efficient
large-cap equity universe with a minimal technical feature set?
\textbf{RQ2:} Does each relational scale contribute independently to
predictive performance, or does performance derive primarily from one
dominant module?
\textbf{RQ3:} Does the architecture transfer across markets with
different microstructural characteristics, or does its effectiveness
depend on specific conditions of sector diversity and market efficiency?

The contributions of this work are as follows:
\begin{enumerate}
\item \textbf{Three-scale parallel relational architecture.} CRML
jointly constructs static group, dynamic group, and directed pairwise
topologies in a single forward pass---a combination absent from prior
work.

\item \textbf{Dual-polarity dynamic hypergraph.} The KNN module builds
separate co-mover and inverse-mover hyperedge channels per forward pass,
allowing the model to learn distinct aggregation functions for momentum
spillover and mean-reversion simultaneously.

\item \textbf{Stock-specific gated fusion.} A gating mechanism computes
dynamic routing weights conditioned on current embeddings, replacing
fixed blending with a learned, per-stock, per-day fusion.

\item \textbf{Rigorous walk-forward evaluation.} Hyperparameter
optimization is confined to folds~1--3; folds~4--7 are withheld
entirely, eliminating the forward-looking bias present in the fixed
train/val/test splits used by MASTER~\cite{li2024master} and
StockMixer~\cite{fan2024stockmixer}.

\item \textbf{Cross-market boundary identification.} A positive result
on the S\&P~100 alongside a null result on the NIFTY~100 identifies the
structural conditions under which multi-scale relational learning does
and does not generate detectable alpha.
\end{enumerate}

The rest of this paper is organized as follows. Section~II surveys
related work. Section~III presents the architecture and evaluation
protocol. Section~IV reports results. Section~V provides direct answers
to the three research questions and concludes.

\section{Related Work}

\subsection{Traditional Machine Learning for Stock Forecasting}

Early quantitative methods applied linear regression and support vector
regression to engineered price features~\cite{gu2020empirical}. These
models are interpretable but structurally limited: they treat each stock
as an independent observation and carry no mechanism for encoding the
co-movement patterns that distinguish equity markets from isolated time
series. Jegadeesh and Titman documented that cross-sectional momentum
persists across multi-month horizons in US markets~\cite{jegadeesh1993returns},
providing the empirical basis for the momentum features used in our
feature set. Grinold and Kahn formalized the IC as the appropriate
measure of cross-sectional predictive power, distinguishing the ranking
problem from scalar return forecasting~\cite{grinold2000active}. Later
work by Gu et al.\ showed that non-linear machine learning methods
substantially outperform linear baselines on US equity cross-sections by
capturing feature interaction effects~\cite{gu2020empirical}. The shared
limitation across all these methods is that they process each stock
independently: two stocks in the same sector, with the same supply chain
exposure, are modeled as if they share no information.

\subsection{Deep Learning for Sequential Return Prediction}

Recurrent architectures replaced hand-crafted temporal features with
learned representations. The Long Short-Term
Memory~\cite{hochreiter1997long} became the default sequential baseline;
its gating mechanism selectively retains trend information across
variable horizons. The Gated Recurrent Unit~\cite{cho2014learning}
achieved similar empirical performance with fewer parameters. Neither
model shares information across stocks during encoding.

The Attention LSTM (ALSTM)~\cite{qin2017dual} added a soft attention
layer over hidden states, improving IC on Qlib benchmarks relative to
vanilla LSTM, but the attention operates within a single stock's
sequence. The State Frequency Memory (SFM)~\cite{zhang2017stock}
decomposed hidden states into frequency components via a Discrete
Fourier Transform formulation, jointly capturing short- and long-term
trading patterns. Both SFM and ALSTM apply independently per stock
with no cross-sectional routing. Transformer
encoders~\cite{vaswani2017attention} extended temporal context through
self-attention, but applied stock-by-stock they still propagate no
cross-sectional information. Cross-stock attention, as used in MASTER,
addresses this at the cost of conflating structurally distinct relational
mechanisms.

\subsection{Graph Neural Networks for Stock Prediction}

Kipf and Welling established GCN as a scalable approach to learning node
representations incorporating neighborhood
information~\cite{kipf2017semi}. The natural application to stocks
treats each asset as a node and encodes known relationships as edges.

RSR~\cite{feng2019temporal} combined LSTM encoders with Temporal Graph
Convolution over a static knowledge graph of sector and supply chain
relations, demonstrating measurable improvement over independent LSTM on
NYSE and NASDAQ data. HATS~\cite{kim2019hats} extended this with
heterogeneous relation types from Wikidata and hierarchical attention to
select among them, outperforming GCN and TGC baselines on S\&P~500 data
from 2013 to 2019. Both share the same structural constraint: the
topology is built once and frozen.

TRA~\cite{lin2021learning} addressed within-universe return
heterogeneity by routing each sample to specialized prediction heads,
showing that stocks with different return dynamics benefit from distinct
model components. DoubleAdapt~\cite{zhao2023doubleadapt} addressed
distribution shift through meta-learned data and model adapters,
reaching strong performance on CSI~300 and CSI~500 under an incremental
learning protocol. Its contribution is adaptation rather than relational
structure: the base encoder still processes stocks independently.

MASTER~\cite{li2024master} bypassed predefined graph construction by
modeling stock correlations through alternating intra- and inter-stock
transformer layers, reaching IC~$=0.059$ on CSI~300. The limitation is
that fully attentive inter-stock aggregation assigns no structural
meaning to the connection weights: sector co-movement and lead-lag
signals receive identical architectural treatment. StockMixer~\cite{fan2024stockmixer}
took the opposing position---a three-stage MLP mixer across indicators,
time steps, and stocks matches MASTER on IC at substantially lower
cost---but the stock-mixing stage is a learned global linear
transformation, not a structured topology.

\subsection{Hypergraph Learning for Stock Prediction}

Encoding a sector of $k$ stocks as a pairwise graph requires
$\binom{k}{2}$ edges and discards the group structure.
HGNN~\cite{feng2019hypergraph} showed that propagating node features
through a hyperedge and back in a single operation preserves
higher-order co-membership that pairwise decomposition cannot recover.

STHAN-SR~\cite{sawhney2021sthan} applied this to stock selection,
constructing hyperedges from sector memberships and corporate relations.
On NASDAQ over six years, STHAN-SR outperformed RSR-I (Sharpe 1.42
vs.\ 1.34). The hyperedges are static, however, and cannot capture
transient cross-sector groupings that arise during macro-driven events.

HIST~\cite{xu2021hist} added a hidden concept module that discovers
latent shared factors, achieving IC~$=0.048$ on CSI~300. Both HIST
modules operate at the group level, neither reconstructs the topology
per forward pass from current embeddings, and neither models asymmetric
pairwise effects.

\subsection{Multi-Scale and Multi-Relational Learning}

R-GCN~\cite{schlichtkrull2018modeling} established that type-specific
transformation matrices for different edge types in knowledge graphs
improve over single-relation GCN, confirming that distinct relational
semantics benefit from distinct aggregation functions. In the financial
domain, STHGNN~\cite{xiang2022temporal} combined corporate relations,
price co-movement, and temporal dependencies in a heterogeneous graph,
outperforming homogeneous baselines---but all relation types remained
encoded as edge attributes within a single pairwise graph, leaving
hyperedge-level group structure unmodeled.

Surveys by Patel et al.~\cite{patel2024gnn} and Zhang
et al.~\cite{zhang2024neural} both identify simultaneous multi-scale
topology construction as an open problem. No prior architecture jointly
builds a static group hypergraph, a dynamic dual-polarity group
hypergraph, and a directed pairwise attention graph, nor fuses them
through a learned per-stock gating mechanism. CRML addresses this gap
directly.

\section{Methodology}

\subsection{Problem Formulation}

Let $\mathcal{U}=\{s_1,s_2,\dots,s_N\}$ denote the stock universe,
where $N=100$ for the S\&P~100 and $N=91$ for the NIFTY~100. At each
trading day $t$, stock $s_i$ has feature matrix
$\mathbf{X}_i\in\mathbb{R}^{T\times F}$, where $T=60$ is the lookback
window and $F=12$ is the feature count. The cross-sectional input
tensor is $\mathbf{X}\in\mathbb{R}^{N\times T\times F}$. A macro
context tensor $\mathbf{M}\in\mathbb{R}^{T\times 7}$ encodes seven
market-level indicators over the same window.

The prediction target is the 5-day simple return:
\begin{equation}
  y_i^t = \frac{C_{t+5}}{C_t} - 1
  \label{eq:target}
\end{equation}
where $C_t$ is the closing price on the last day of the input window.
A one-day safety gap is enforced: the first return contributing to
$y_i^t$ is $C_{t+1}/C_t$, and the five consecutive log returns
telescope to $\ln(C_{t+5}/C_t)$, which exponentiates to
$(C_{t+5}/C_t)-1$. IC and Rank\,IC metrics are scale-invariant and
computed on unscaled returns; the reported IC of 0.047 is unaffected
by the loss-level scaling described in Section~\ref{sec:loss}.

The model produces a magnitude prediction $\hat{y}_i\in\mathbb{R}$ and
a direction prediction $\hat{d}_i\in[0,1]$. The primary metric is the
cross-sectional IC~\cite{grinold2000active} via Spearman rank
correlation~\cite{spearman1904}:
\begin{equation}
  \mathrm{IC}_t = \mathrm{SpearmanCorr}\!\left(
                   \hat{\mathbf{y}}^t,\,\mathbf{y}^t\right)
  \label{eq:ic}
\end{equation}
Mean IC across all test days measures average cross-sectional ranking
ability independent of market direction~\cite{grinold2000active}.

\subsection{Feature Engineering}

\textbf{Stock-level features.} Each stock has 12 technical features per
day (Table~\ref{tab:features}). Momentum features~\cite{jegadeesh1993returns}
capture return persistence at 5-, 10-, and 20-day horizons. The RSI is
per Wilder~\cite{wilder1978rsi} and normalized to $[0,1]$. MACD follows
Appel~\cite{appel2005technical}, Bollinger Band position follows
Bollinger~\cite{bollinger2001bands}, and the normalized ATR is also
from Wilder~\cite{wilder1978rsi}. All features use only past data, with
a one-day shift on all rolling statistics to eliminate look-ahead bias.
Feature values are cross-sectionally z-scored per day---a standard
technique for removing market beta from factor
signals~\cite{kakushadze2020101}. Missing values are forward-filled up
to five days; longer gaps are zero-filled.

\begin{table}[!t]
\renewcommand{\arraystretch}{1.25}
\caption{Stock-Level Features (Cross-Sectionally Z-Scored Daily)}
\label{tab:features}
\centering
\resizebox{\columnwidth}{!}{%
\begin{tabular}{llcl}
\toprule
\textbf{Feature} & \textbf{Definition} & \textbf{Window} & \textbf{Type}\\
\midrule
\texttt{log\_ret}      & $\ln(c_t/c_{t-1})$                          & 1\,d     & Price\\
\texttt{volatility}    & Rolling std of log returns                  & 20\,d    & Risk\\
\texttt{rsi\_norm}     & Normalized RSI $\in[0,1]$~\cite{wilder1978rsi} & 14\,d & Momentum\\
\texttt{mom\_5}        & $(c_t/c_{t-5})-1$~\cite{jegadeesh1993returns}  & 5\,d  & Momentum\\
\texttt{mom\_10}       & $(c_t/c_{t-10})-1$                          & 10\,d    & Momentum\\
\texttt{mom\_20}       & $(c_t/c_{t-20})-1$                          & 20\,d    & Momentum\\
\texttt{vol\_surprise} & $v_t/\mathrm{EWM}(v,\mathrm{span}{=}20)$   & 20\,d    & Volume\\
\texttt{ma\_dist\_20}  & $(c-\mathrm{MA}_{20})/\mathrm{MA}_{20}$    & 20\,d    & Trend\\
\texttt{ma\_dist\_50}  & $(c-\mathrm{MA}_{50})/\mathrm{MA}_{50}$    & 50\,d    & Trend\\
\texttt{macd}          & $\mathrm{EMA}_{12}-\mathrm{EMA}_{26}$~\cite{appel2005technical} & 12/26\,d & Trend\\
\texttt{bb\_position}  & $(c-L)/(U-L)$~\cite{bollinger2001bands}     & 20\,d    & Volatility\\
\texttt{atr\_norm}     & $\mathrm{ATR}_{14}/c$~\cite{wilder1978rsi}  & 14\,d    & Risk\\
\bottomrule
\end{tabular}%
}
\end{table}

\textbf{Macro features.} Seven market-level indicators per market
capture systemic macroeconomic risk factors~\cite{chen1986economic}.
For the US: CBOE VIX~\cite{whaley2000vix}, S\&P~500 index, 10-year
Treasury yield, US Dollar Index, gold futures, crude oil futures, and
Nasdaq~100. For India: India~VIX, NIFTY~50, Bank~NIFTY, USD/INR
exchange rate, gold futures, crude oil futures, and the 10-year US
Treasury yield. Table~\ref{tab:macro} lists each indicator with its
economic rationale. Macro series are z-scored over a rolling 60-day
window with the same one-day shift.

\begin{table}[!t]
\renewcommand{\arraystretch}{1.25}
\caption{Macro-Economic Context Indicators (7 per Market)}
\label{tab:macro}
\centering
\resizebox{\columnwidth}{!}{%
\begin{tabular}{llll}
\toprule
\textbf{Indicator} & \textbf{US Ticker} & \textbf{India Ticker}
  & \textbf{Economic Rationale}\\
\midrule
Equity volatility index
  & \texttt{\^{}VIX} & \texttt{\^{}INDIAVIX}
  & Implied market fear; spikes precede broad drawdowns~\cite{whaley2000vix}\\
Broad market index
  & \texttt{\^{}GSPC} & \texttt{\^{}NSEI}
  & Overall market beta and trend regime\\
Sector benchmark
  & \texttt{\^{}NDX} & \texttt{\^{}NSEBANK}
  & Growth/tech proxy (US); financial-sector proxy (India)\\
Sovereign yield
  & \texttt{\^{}TNX} & \texttt{\^{}TNX}
  & Risk-free rate; rising yields compress equity multiples\\
Currency index
  & \texttt{DX-Y.NYB} & \texttt{USDINR=X}
  & Dollar strength affects multinational earnings;
    INR/USD reflects FII flow pressure\\
Gold futures
  & \texttt{GC=F} & \texttt{GC=F}
  & Safe-haven demand; inversely correlated with risk appetite\\
Crude oil futures
  & \texttt{CL=F} & \texttt{CL=F}
  & Input cost for energy-intensive sectors; inflation signal\\
\bottomrule
\end{tabular}%
}
\end{table}

\textbf{Target construction.} The target $y_i^t=(C_{t+5}/C_t)-1$ is
computed on close-to-close prices with the one-day gap. All ranking
metrics (IC, Rank\,IC, Sharpe, hit rate) use unscaled decimal returns.
RMSE and MAE are in natural return units.

\begin{figure*}[!t]
\centering
\input{architecture.tex}
\caption{The CRML architecture. Each stock's 60-day feature tensor
$\mathbf{X}\in\mathbb{R}^{N\times60\times12}$ is compressed by a shared
GRU encoder into temporal embeddings $\mathbf{H}\in\mathbb{R}^{N\times d}$.
Three parallel modules construct a static GICS sector hypergraph
($\mathbf{H}_\mathrm{sec}$), a dynamic dual-polarity KNN hypergraph rebuilt
each forward pass ($\mathbf{H}_\mathrm{knn}$), and a sparse GATv2 pairwise
attention graph ($\mathbf{H}_\mathrm{gat}$). A gating layer produces
per-stock softmax weights and computes the fused embedding $\mathbf{h}_i^*$.
Two independent heads optimize magnitude (regression) and direction
(classification) via a four-component loss.}
\label{fig:architecture}
\end{figure*}

\subsection{GRU Temporal Encoder}

A single-layer GRU~\cite{cho2014learning} compresses each stock's
feature sequence into a final hidden state $\mathbf{h}_i\in\mathbb{R}^{d}$,
where $d$ is set by HPO (Table~\ref{tab:hpo}). At $T=60$ and $F=12$,
sequences are short enough that GRU's streamlined gating captures the
relevant temporal dependencies without the full cell state of an
LSTM~\cite{chung2014empirical}. Deploying independent self-attention
encoders per stock would add an $\mathcal{O}(T^2)$ bottleneck for each
of $N$ stocks per forward pass~\cite{vaswani2017attention}---unjustified
at this sequence length. All $N$ hidden states are stacked into
$\mathbf{H}\in\mathbb{R}^{N\times d}$.

\subsection{Three Parallel Relational Modules}

Each module takes $\mathbf{H}$ as input, constructs a distinct relational
structure, and returns an updated embedding matrix of shape
$\mathbb{R}^{N\times d}$.

\subsubsection{Static Sector Hypergraph}

Stocks are grouped by their Global Industry Classification Standard
sector~\cite{msci2023gics}. Each sector forms a hyperedge connecting
all member stocks, so the entire sector responds jointly to shared
fundamental drivers: regulatory shifts, commodity input costs, interest
rate sensitivity. Encoding a sector of $k$ stocks as a pairwise clique
requires $\binom{k}{2}$ edges; a single hyperedge replaces all of them.

The incidence matrix $\mathbf{H}_s\in\{0,1\}^{N\times S}$ has
$\mathbf{H}_s[i,j]=1$ if stock $i$ belongs to sector $j$. Sector labels
were obtained from Yahoo Finance at the time of data download and
reflect current GICS classifications; historical reclassifications are
not accounted for. Message passing follows HGNN~\cite{feng2019hypergraph}:
\begin{equation}
  \mathbf{X}' = \mathbf{D}_v^{-1/2}\,\mathbf{H}_s\,\mathbf{W}\,
                \mathbf{D}_e^{-1}\,\mathbf{H}_s^\top\,
                \mathbf{D}_v^{-1/2}\,\mathbf{X}\,\boldsymbol{\Theta}
  \label{eq:hgnn_static}
\end{equation}
where $\mathbf{D}_v$ and $\mathbf{D}_e$ are diagonal degree matrices,
$\mathbf{W}=\mathbf{I}_S$, and
$\boldsymbol{\Theta}\in\mathbb{R}^{d\times d}$ is a learnable
projection. Two layers are stacked with ReLU and dropout ($p=0.30$)
between them. The topology is frozen; sector assignments do not change
per forward pass.

\subsubsection{Dynamic Dual-Polarity KNN Hypergraph}

Empirical correlation structures shift with market conditions. Stocks
that co-move during a macro-driven selloff may decorrelate entirely
during a sector rotation days later. To capture this, the KNN
hypergraph is rebuilt from scratch on every forward
pass~\cite{wang2019dynamic}.

GRU embeddings are first passed through a learned linear layer and
$\ell_2$-normalized, then used to compute pairwise cosine similarity:
\begin{equation}
  \tilde{\mathbf{h}}_i = \frac{\mathbf{W}_{\mathrm{proj}}\mathbf{h}_i}
    {\|\mathbf{W}_{\mathrm{proj}}\mathbf{h}_i\|},\qquad
  \mathbf{S}[i,j] = \tilde{\mathbf{h}}_i^\top\tilde{\mathbf{h}}_j
  \label{eq:cosine_sim}
\end{equation}
where $\mathbf{W}_{\mathrm{proj}}$ is learnable. A separate projection
lets the model tune the similarity metric for the relational task
without disturbing the temporal encoding.

For each stock $i$, two hyperedge sets are built. A \textit{positive
hyperedge} connects $i$ to its $k$ most similar stocks, capturing
momentum spillover and co-movement~\cite{jegadeesh1993returns}. A
\textit{negative hyperedge} connects $i$ to its $k$ least similar
stocks, capturing mean-reversion and divergence. Negative hyperedge
indices are offset by $N$ so the HGNN can learn a distinct aggregation
function for each polarity. Every stock also has a self-loop. The
dual-polarity incidence matrix $\mathbf{H}_k\in\{0,1\}^{N\times3N}$
is rebuilt every forward pass; message passing follows
Eq.~(\ref{eq:hgnn_static}) with $\mathbf{H}_k$ in place of
$\mathbf{H}_s$. The HPO-selected value $k=20$ yields $3N$ hyperedges
per pass ($N$ positive, $N$ negative, $N$ self-loops). Values below~5
produce degenerate isolated cliques; values above~20 over-smooth the
co-mover hyperedges toward a fully connected state and collapse the
inverse-mover hyperedges toward near-random pairings.

\subsubsection{Pairwise GATv2}

The hypergraph modules perform symmetric group aggregation, which
cannot represent directed influence: a large-cap bellwether predicts
its supply chain's returns, but the reverse does not hold in general.
The GATv2 module models these directed, asymmetric dependencies.

From $\mathbf{S}$, the top $K=10$ edges per node form a sparse
directed graph. GATv2Conv~\cite{brody2022attentive} with two attention
heads is applied:
\begin{equation}
  \alpha_{ij} = \mathrm{softmax}_j\!\left(
    \mathbf{a}^\top\,\mathrm{LeakyReLU}\!\left(
      \mathbf{W}[\mathbf{h}_i\parallel\mathbf{h}_j]\right)\right)
  \label{eq:gatv2}
\end{equation}
GATv2 is used instead of GAT~v1~\cite{velickovic2018graph} because
v1's static attention cannot distinguish neighbors based on the query
node's current state~\cite{brody2022attentive}. Two GATv2 layers are
applied; the first concatenates heads to $2d$ and the second averages
back to $d$. The output $\mathbf{H}_{\mathrm{gat}}\in\mathbb{R}^{N\times d}$
encodes directed pairwise attention weights.

\subsection{Gated Fusion}

The three outputs are combined through an input-conditioned gate rather
than fixed concatenation, which would impose a static blending ratio
across all stocks and regimes. For each stock $i$:
\begin{align}
  \mathbf{z}_i &= [\mathbf{h}_{\mathrm{sec},i}\parallel
                   \mathbf{h}_{\mathrm{knn},i}\parallel
                   \mathbf{h}_{\mathrm{gat},i}]\in\mathbb{R}^{3d}
                   \label{eq:concat}\\
  \mathbf{w}_i &= \mathrm{softmax}(\mathbf{W}_g\mathbf{z}_i +
                   \mathbf{b}_g)\in\mathbb{R}^3
                   \label{eq:gate_weights}\\
  \mathbf{h}_i^* &= \textstyle\sum_c w_{i,c}\,\mathbf{h}_{c,i}
                   \in\mathbb{R}^{d}
                   \label{eq:fuse}
\end{align}
where $\mathbf{W}_g\in\mathbb{R}^{3\times3d}$ and
$\mathbf{b}_g\in\mathbb{R}^3$ are shared across all stocks. Gate
weights are recomputed on every forward pass, making them stock-specific
and state-dependent. Over the S\&P~100 evaluation period, the optimizer
converges to a near-uniform equilibrium
($\mathbf{w}_i\approx[0.333,0.333,0.333]$), discussed in
Section~\ref{sec:ablation}.

\subsection{Prediction Heads}

The fused embedding $\mathbf{h}_i^*$ feeds two independent networks:
\begin{itemize}
  \item \textbf{Magnitude head:} $\mathbb{R}^d\to\mathbb{R}^{64}\to
    \mathrm{ReLU}\to\mathrm{Dropout}(p)\to\mathbb{R}^1$,
    producing $\hat{y}_i$.
  \item \textbf{Direction head:} same structure with a sigmoid output,
    producing $\hat{d}_i\in[0,1]$.
\end{itemize}
The direction head is an auxiliary regularizer~\cite{caruana1997multitask}.
Requiring $\mathbf{h}_i^*$ to solve a simultaneous binary
classification problem constrains the embedding to encode directional
boundaries, stabilizing gradient flow and reducing sensitivity to
return outliers. All reported ranking metrics use only $\hat{y}_i$.

\subsection{Multi-Component Loss Function}
\label{sec:loss}

\begin{equation}
  \mathcal{L} = w_{\mathrm{IC}}\mathcal{L}_{\mathrm{IC}}
              + w_{\mathrm{MSE}}\mathcal{L}_{\mathrm{MSE}}
              + w_{\mathrm{Rank}}\mathcal{L}_{\mathrm{Rank}}
              + w_{\mathrm{BCE}}\mathcal{L}_{\mathrm{BCE}}
  \label{eq:total_loss}
\end{equation}

\textbf{IC loss} ($w_{\mathrm{IC}}=0.31$): negative Pearson
correlation~\cite{pearson1895notes} between predicted and actual returns
in the batch---a differentiable proxy directly aligned with the IC
evaluation metric:
\begin{equation}
  \mathcal{L}_{\mathrm{IC}} = -\frac{
    \sum_i(\hat{y}_i-\bar{\hat{y}})(y_i-\bar{y})}{
    \sqrt{\sum_i(\hat{y}_i-\bar{\hat{y}})^2\cdot
          \sum_i(y_i-\bar{y})^2}}
  \label{eq:loss_ic}
\end{equation}

\textbf{MSE loss} ($w_{\mathrm{MSE}}=0.49$): mean squared error between
$\hat{y}_i$ and the target return scaled by~100. This scaling is applied
only inside the MSE component to condition gradients during early
training and does not affect the IC, Rank, or BCE components. Without
MSE, pure IC optimization collapses into a rank-only solution that
ignores return magnitude. The dominance of $w_{\mathrm{MSE}}$ over
$w_{\mathrm{IC}}$ also prevents gradient explosions when the denominator
of Eq.~(\ref{eq:loss_ic}) approaches zero in early training.

\textbf{Pairwise ranking loss} ($w_{\mathrm{Rank}}=0.10$): a
margin-based penalty from the learning-to-rank
paradigm~\cite{burges2005learning}. For all pairs $(i,j)$ with
$y_i>y_j$, it penalizes $\hat{y}_i\le\hat{y}_j$:
\begin{equation}
  \mathcal{L}_{\mathrm{Rank}} = \frac{1}{|\mathcal{P}|}
    \sum_{(i,j)\in\mathcal{P}}
    \max\!\left(0,\,-(\hat{y}_i-\hat{y}_j)
    \cdot\mathrm{sign}(y_i-y_j)+m\right)
  \label{eq:loss_rank}
\end{equation}
where $\mathcal{P}$ is all pairs with $y_i\ne y_j$ and $m$ is the
margin.

\textbf{Binary cross-entropy} ($w_{\mathrm{BCE}}=0.10$): cross-entropy
between $\hat{d}_i$ and $\mathbf{1}[y_i>0]$, acting as a representation
regularizer for the direction head. The empirical consequences of
replacing this objective with the standard Qlib protocol are quantified
in Section~\ref{sec:protocol}.

\subsection{Walk-Forward Evaluation Protocol}

Standard evaluation in graph-based stock prediction frequently inflates
IC through hyperparameter leakage~\cite{lopezdeprado2018advances}. We
eliminate this with strict chronological separation.

\textbf{Expanding folds.} Seven folds are constructed. The initial
training window is 500~trading days; each fold advances the test window
by 125~days ($\approx$six months). A 60-day embargo buffer separates
training from test to prevent sequence overlap.

\begin{figure}[!t]
\centering
\resizebox{\columnwidth}{!}{\input{timeline}}
\caption{Expanding walk-forward evaluation protocol. A 60-day embargo
buffer separates all training and test windows. Folds~1--3~(blue) are
reserved for hyperparameter optimization; folds~4--7~(teal) are
withheld until HPO concludes.}
\label{fig:walkforward}
\end{figure}

\textbf{HPO/test separation.} Folds~1--3 are reserved exclusively for
Optuna TPE hyperparameter optimization (50~trials). Folds~4--7 are
embargoed until HPO concludes. No parametric decision uses fold~4--7
data. This eliminates the forward-looking bias in the fixed splits used
by MASTER and StockMixer, where hyperparameter selection implicitly
optimizes against the reported test period.

\textbf{Per-day IC.} We compute IC per trading day (Eq.~\ref{eq:ic})
and report the mean and standard deviation over folds~4--7. Pooled
IC---computed by stacking all predictions across days---conflates
cross-sectional ranking ability with time-series regime
effects~\cite{grinold2000active} and is not reported.

\textbf{Bootstrap CIs.} 95\,\% confidence intervals are from
10\,000 resamples with replacement~\cite{efron1994introduction}.
The US CI of $[0.016,\,0.077]$ lies entirely above zero.

\subsection{Expressivity Analysis}

Standard message-passing neural networks are bounded by the 1-WL graph
isomorphism test~\cite{xu2019how}: neighborhoods that 1-WL cannot
distinguish receive identical embeddings regardless of network depth.
Hypergraph networks on hyperedges of size $>2$ are strictly more
expressive because group-level aggregation distinguishes configurations
that pairwise message-passing conflates~\cite{feng2019hypergraph}.

\smallskip
\noindent\textbf{Proposition~1 (Expressivity Separation).}
\textit{Let $\mathcal{F}_{\mathrm{single}}^A$ be the function class of
single-scale variant
$A\in\{\mathrm{Static\text{-}SecHG},\,\mathrm{KNN\text{-}HG},\,
\mathrm{Pairwise\text{-}GAT}\}$, and let $\mathcal{F}_{\mathrm{multi}}$
be the function class of CRML. Then
$\mathcal{F}_{\mathrm{single}}^A\subsetneq\mathcal{F}_{\mathrm{multi}}$
for each $A$.}

\smallskip
\noindent\textit{Justification.} Subsumption holds because CRML reduces
to any single-scale variant by routing gate weights to a one-hot vector.
For strict inclusion, consider six stocks across two sectors where a
macro shock pushes cross-sector cosine similarity above intra-sector
similarity. A ranking function requiring sector-level group aggregation,
dynamic cross-sector similarity, and directed pairwise asymmetry
simultaneously cannot be realized by any single variant: Static-SecHG
cannot capture transient cross-sector grouping; KNN-HG cannot recover
permanent sector structure when cross-sector similarity dominates;
Pairwise-GAT cannot distinguish a group hyperedge from a star graph
when intra-sector embeddings are equal. CRML computes all three jointly.
Empirical support is in Section~\ref{sec:ablation}.

\smallskip
\noindent\textbf{Remark.} The empirical gate equilibrium (0.333 per
module) is an optimization outcome, not an expressivity failure. The
three scales contribute roughly equal average information on S\&P~100
large caps; the optimizer finds no incentive to break symmetry.
Regime-dependent specialization may emerge under longer training windows
or richer feature sets.

\subsection{Complexity Analysis}

Total forward-pass complexity is $\mathcal{O}(NTd^2+N^2d+Nd^2+NKd)$.
The bottleneck is the GRU encoder at $\mathcal{O}(NTd^2)$; the
relational modules together contribute $\mathcal{O}(Nd^2)$, cheaper by
a factor of $T=60$. For $N=100$, $d=256$, a single-day forward pass
completes in milliseconds on a Tesla~T4.
MASTER~\cite{li2024master} computes full attention over all
$N\!\times\!T$ tokens at $\mathcal{O}(N^2T^2d)$; CRML's sparse routing
reduces the spatial term to $\mathcal{O}(N^2)$. Total parameters remain
under~500K.

% \FloatBarrier
\input{pseudocode}
% \FloatBarrier

\section{Experiments}

\subsection{Datasets and Implementation Details}

\textbf{Datasets.} We use two equity universes. The US experiment uses
the S\&P~100. The India experiment uses the NIFTY~100; 91~of
approximately 100~constituents satisfied the minimum 1\,500-day
continuous history requirement (Tata Motors~DVR, Jio Financial, Zomato,
LIC, Lodha, IRFC, IRCTC, Mankind Pharma, and Max Healthcare were
excluded for insufficient history or download failure). Both datasets
span January~2015 to early~2026, sourced from the Yahoo Finance
API~\cite{yang2020qlib}. Universe constituents reflect index composition
at the download date; stocks delisted or removed between 2015 and 2026
are absent, introducing potential survivorship
bias~\cite{brown1992survivorship} noted as a limitation.
Table~\ref{tab:datasets} summarizes dataset properties.

\begin{table}[!t]
\renewcommand{\arraystretch}{1.25}
\caption{Dataset Characteristics and Evaluation Horizons}
\label{tab:datasets}
\centering
\resizebox{\columnwidth}{!}{%
\begin{tabular}{lll}
\toprule
\textbf{Property} & \textbf{US Market} & \textbf{India Market}\\
\midrule
Universe       & S\&P~100 ($N=100$)  & NIFTY~100 ($N=91$ surviving)\\
Data Source    & Yahoo Finance API   & Yahoo Finance API\\
Date Range     & 2015-01-01 to 2026  & 2015-01-01 to 2026\\
Asset Features & 12 technical        & 12 technical\\
Macro Context  & VIX, SPX, TNX, DXY, & INDIAVIX, NSEI, NSEBANK,\\
               & Gold, Oil, NDX      & USDINR, Gold, Oil, TNX\\
Target Horizon & 5-day return        & 5-day return\\
Lookback ($T$) & 60 trading days     & 60 trading days\\
\bottomrule
\end{tabular}%
}
\end{table}

\textbf{Hyperparameters.} US hyperparameters were selected via
50-trial Optuna TPE~\cite{akiba2019optuna} on folds~1--3. India HPO
converged to a degenerate state---the top~10 trials produced identical
configurations---and provided no improvement over the US configuration.
Both are reported in the cross-market analysis.
Table~\ref{tab:hpo} lists finalized settings.

\begin{table}[!t]
\renewcommand{\arraystretch}{1.25}
\caption{Post-HPO Hyperparameter Configurations}
\label{tab:hpo}
\centering
\resizebox{\columnwidth}{!}{%
\begin{tabular}{lcc}
\toprule
\textbf{Hyperparameter} & \textbf{US (S\&P~100)} & \textbf{India (NIFTY~100)}\\
\midrule
Hidden Dimension ($d$)   & 256                 & 128\\
Learning Rate            & $5\times10^{-5}$    & $1\times10^{-4}$\\
Dropout ($p$)            & 0.30                & 0.30\\
KNN Neighbors ($k$)      & 20                  & 20\\
GAT Top Edges ($K$)      & 10                  & 10\\
Loss ($w_\mathrm{IC}/w_\mathrm{MSE}/w_\mathrm{Rank}/w_\mathrm{BCE}$)
                         & 0.31/0.49/0.10/0.10 & 0.50/0.30/0.10/0.10\\
Epochs per Fold          & 20                  & 25\\
Sequence Length ($T$)    & 60                  & 60\\
Batch Size               & 64                  & 64\\
Gradient Clipping        & 1.1                 & 1.1\\
Weight Decay             & $4\times10^{-6}$    & $4\times10^{-6}$\\
Optimizer                & \multicolumn{2}{c}{AdamW~\cite{loshchilov2017decoupled}
                           + Cosine Annealing}\\
\bottomrule
\end{tabular}%
}
\end{table}

\textbf{Implementation.} All experiments use Python~3 with
PyTorch~\cite{paszke2019pytorch} and PyTorch
Geometric~\cite{fey2019fast} on a single NVIDIA Tesla~T4 GPU
(16\,GB VRAM). Manual random seeds and
\texttt{torch.backends.cudnn.deterministic=True} ensure reproducibility.

\subsection{Main Results --- US Market (S\&P~100)}

Table~\ref{tab:main_results} reports out-of-sample performance on
folds~4--7. CRML achieves mean daily IC~$=0.047$ with 95\,\% bootstrap
CI~$[0.016,\,0.077]$. The lower bound is strictly positive, confirming
a statistically significant cross-sectional ranking signal. A long-short
decile portfolio yields Sharpe~$=1.263$ and hit rate~$=56.98\,\%$.
Ridge regression reaches IC~$=0.020$ under identical features and
protocol; the gap (0.027 IC points) represents the contribution of
CRML's relational routing over linear feature aggregation.

\begin{table}[!t]
\renewcommand{\arraystretch}{1.25}
\caption{Main Results on S\&P~100 (Held-Out Test Folds~4--7)}
\label{tab:main_results}
\centering
\resizebox{\columnwidth}{!}{%
\begin{tabular}{lccccc}
\toprule
\textbf{Model}
  & \textbf{IC (Mean$\pm$Std)}
  & \textbf{95\,\% CI}
  & \textbf{Sharpe}
  & \textbf{Hit Rate}
  & \textbf{L/S Spread}\\
\midrule
\textbf{CRML (Proposed)}
  & $\mathbf{0.047\pm0.033}$
  & $\mathbf{[0.016,\,0.077]}$
  & $\mathbf{1.263}$
  & $\mathbf{56.98\,\%}$
  & $\sim\!\mathbf{0.005}$\\
Linear Baseline (Ridge)
  & $\sim\!0.020$ & -- & -- & -- & --\\
\bottomrule
\end{tabular}%
}
\end{table}

\begin{figure}[!t]
\centering
\includegraphics[width=\columnwidth]{graphs/fig_us_ic_bar.png}
\caption{Mean daily IC per fold, S\&P~100. Blue~=~HPO folds~1--3
(excluded from primary results); orange~=~test folds~4--7. Dashed
lines mark the ridge baseline (IC~$=0.020$) and a reference threshold
(IC~$=0.050$). Fold~6 (IC~$=0.004$) is near zero; Fold~7 (IC~$=0.089$)
is the strongest. The fold-level spread reflects the non-stationarity
of cross-sectional predictability documented across
models~\cite{gu2020empirical}.}
\label{fig:us_ic_bar}
\end{figure}

\begin{figure}[!t]
\centering
\includegraphics[width=\columnwidth]{graphs/fig_us_daily_ic_ts.png}
\caption{Daily IC time series, test folds~4--7, S\&P~100. Each line
is one 125-trading-day fold. Fold~6 sustains near-zero mean IC
throughout; Fold~7 is consistently positive. The contrast between
folds is consistent with macro-regime sensitivity.}
\label{fig:us_daily_ic_ts}
\end{figure}

\begin{figure}[!t]
\centering
\includegraphics[width=\columnwidth]{graphs/fig_us_bootstrap_ic_hist.png}
\caption{Bootstrap IC distribution, S\&P~100 (10\,000 resamples).
Vertical lines mark the sample mean~(0.047) and CI bounds
$[0.016,\,0.077]$. Both bounds lie above zero.}
\label{fig:us_bootstrap_ic_hist}
\end{figure}

Fold-level IC varies substantially (Table~\ref{tab:perfoldc}).
Fold~7 (IC~$=0.089$, Sharpe~$=2.417$) and Fold~5 (IC~$=0.065$,
Sharpe~$=2.241$) are strongest; Fold~6 (IC~$=0.004$) is the weakest.
This spread is typical of rolling-window equity prediction across
models~\cite{gu2020empirical}: no architecture extracts uniform IC
across all regimes.

\begin{table}[!t]
\renewcommand{\arraystretch}{1.25}
\caption{Per-Fold IC Breakdown (S\&P~100)}
\label{tab:perfoldc}
\centering
\resizebox{\columnwidth}{!}{%
\begin{tabular}{clccl}
\toprule
\textbf{Fold} & \textbf{Role} & \textbf{IC} & \textbf{Sharpe} & \textbf{Notes}\\
\midrule
1 & HPO & 0.041 & 1.402 & Excluded from reported results\\
2 & HPO & 0.026 & 0.903 & Excluded from reported results\\
3 & HPO & 0.019 & 0.360 & Excluded from reported results\\
\midrule
4 & \textbf{Test} & 0.028 & 0.195 & Out-of-sample\\
5 & \textbf{Test} & 0.065 & 2.241 & Out-of-sample\\
6 & \textbf{Test} & 0.004 & 0.197 & Near-zero; regime-sensitive\\
7 & \textbf{Test} & 0.089 & 2.417 & Strongest fold\\
\bottomrule
\end{tabular}%
}
\end{table}

\subsection{Ablation Study}
\label{sec:ablation}

Eight model variants are evaluated under identical post-HPO
hyperparameters and compared against the full architecture using
one-sided paired t-tests ($\alpha=0.10$) across folds~4--7. Cohen's
$d$ is reported alongside $p$-values because the low fold count ($n=4$)
limits statistical power. Table~\ref{tab:ablation} reports results.

\begin{table*}[!t]
\renewcommand{\arraystretch}{1.25}
\caption{Ablation Study, S\&P~100 (Test Folds~4--7).
         $^\dagger$The ablation full-model IC of 0.043 differs from the
         main result IC of 0.047 because it is from a separate training
         run; the difference lies within fold-level variance and does
         not affect relative comparisons.}
\label{tab:ablation}
\centering
\begin{tabular}{lccccccr}
\toprule
\textbf{Variant}
  & \textbf{IC (Mean$\pm$Std)}
  & $\boldsymbol{\Delta}$\textbf{IC}
  & \textbf{$p$}
  & \textbf{Cohen's $d$}
  & \textbf{Sig.\ ($\alpha{=}0.10$)}
  & \textbf{Sharpe}
  & \textbf{Hit Rate}\\
\midrule
\textbf{Full Model}$^\dagger$
  & $\mathbf{0.043\pm0.029}$ & --   & --    & --    & --  & $\mathbf{1.238}$ & $\mathbf{57.0\,\%}$\\
\midrule
No Static Sector HG
  & $0.031\pm0.025$ & $-0.012$ & 0.053 & 1.327 & Yes & 0.899 & 57.0\,\%\\
No Dynamic KNN HG
  & $0.037\pm0.030$ & $-0.006$ & 0.014 & 2.323 & Yes & 1.104 & 57.0\,\%\\
No Pairwise GAT
  & $0.039\pm0.037$ & $-0.005$ & 0.198 & 0.570 & No  & 1.091 & 57.0\,\%\\
\midrule
Static Sector HG Only
  & $0.038\pm0.040$ & $-0.005$ & 0.336 & 0.270 & No  & 1.039 & 57.0\,\%\\
Dynamic KNN HG Only
  & $0.029\pm0.022$ & $-0.014$ & 0.038 & 1.539 & Yes & 0.857 & 57.0\,\%\\
Pairwise GAT Only
  & $0.033\pm0.022$ & $-0.010$ & 0.082 & 1.059 & Yes & 0.934 & 57.0\,\%\\
GRU Only (No Graph)
  & $0.035\pm0.027$ & $-0.008$ & 0.079 & 1.081 & Yes & 1.097 & 57.0\,\%\\
\bottomrule
\end{tabular}
\end{table*}

\begin{figure}[!t]
\centering
\includegraphics[width=\columnwidth]{graphs/fig_us_ablation_bar.png}
\caption{Mean IC across ablation variants, S\&P~100. Dashed line marks
the full model (0.043). Every variant falls below it. KNN HG alone
(0.029) falls below GRU-only (0.035), confirming that dynamic
similarity clustering requires stable sector structure to add value.}
\label{fig:us_ablation_bar}
\end{figure}

\begin{figure}[!t]
\centering
\includegraphics[width=\columnwidth]{graphs/fig_us_sharpe_bar.png}
\caption{Sharpe ratio across ablation variants. GAT removal misses
IC significance ($p=0.198$) but drops Sharpe from 1.238 to 1.091
(13.5\,\% degradation). Directed bilateral signals improve portfolio
efficiency even when their raw IC contribution is small.}
\label{fig:us_sharpe_bar}
\end{figure}

\begin{figure*}[!t]
\centering
\includegraphics[width=\textwidth]{graphs/fig_us_ablation_heatmap.png}
\caption{IC heatmap: 8 ablation variants (rows) $\times$ 4 test folds
(columns), S\&P~100. Color: red (low) to green (high IC). Near-zero IC
across all variants in Fold~6 points to a macro regime unfavorable to
cross-sectional prediction---not an architectural problem. The full
model row is consistently the darkest.}
\label{fig:us_ablation_heatmap}
\end{figure*}

\begin{figure}[!t]
\centering
\includegraphics[width=\columnwidth]{graphs/fig_fusion_gate_weights.png}
\caption{Fusion gate weights averaged over all stocks and test days,
S\&P~100. The near-uniform distribution ($\approx\!0.333$ per module)
is an optimization equilibrium. All three scales carry complementary
information---removing any one module degrades performance.}
\label{fig:fusion_gate_weights}
\end{figure}

Every variant that removes or isolates a module produces lower mean IC,
so the tri-scale combination outperforms all single-scale and dual-scale
alternatives tested. The Static Sector Hypergraph is the primary anchor:
its removal causes the largest absolute drop ($\Delta\mathrm{IC}=-0.012$,
$p=0.053$), and it is the only module that exceeds the GRU-only baseline
in isolation (0.038 vs.\ 0.035). This is consistent with the established
dominance of sector co-movement in large-cap US equity
returns~\cite{jegadeesh1993returns}.

The Dynamic KNN Hypergraph exhibits a structural duality. Removing it
produces the largest effect size ($p=0.014$, Cohen's $d=2.323$), but
using it alone falls below the GRU-only baseline (0.029 vs.\ 0.035).
Dynamic similarity is a strong auxiliary signal when sector structure
provides a stable anchor; without that anchor, it introduces noise.

GAT removal does not reach IC significance ($p=0.198$), but its
inclusion raises Sharpe from 1.091 to 1.238 (13.5\,\%). Directed
bilateral signals improve portfolio efficiency even when their IC
contribution is statistically marginal.

The gate equilibrium (0.333 per module) is consistent with the ablation
finding: all three scales carry complementary information, so the
optimizer has no incentive to silence any of them. Regime-dependent
gate specialization via explicit symmetry-breaking remains future work.

\subsection{Training Protocol Analysis}
\label{sec:protocol}

MASTER, StockMixer, and HIST all use the standard Qlib training
protocol: MSE-only loss with \textit{CSZScoreNorm} label normalization
and \textit{DropExtremeLabel} preprocessing~\cite{yang2020qlib}. Four
variants were tested to determine whether adopting this protocol would
improve CRML's IC.

\begin{table}[!t]
\renewcommand{\arraystretch}{1.25}
\caption{Training Protocol Comparison (S\&P~100, Test Folds~4--7).
         Variants~B--D were not HPO-tuned; the scale of IC collapse
         far exceeds what hyperparameter mismatch alone would explain.}
\label{tab:training_method}
\centering
\resizebox{\columnwidth}{!}{%
\begin{tabular}{cllccc}
\toprule
\textbf{Var.} & \textbf{Loss} & \textbf{Label Processing}
              & \textbf{IC} & \textbf{Sharpe} & \textbf{Hit Rate}\\
\midrule
A & Multi-objective (ours) & Raw $\times100$            & 0.043 & 1.218 & 57.0\,\%\\
B & Multi-objective        & CSZScoreNorm               & 0.032 & 1.190 & 57.0\,\%\\
C & MSE only (Qlib)        & CSZScoreNorm + DropExtreme & 0.007 & 0.374 & 46.7\,\%\\
D & MSE + 0.1$\times$Rank  & CSZScoreNorm + DropExtreme & 0.006 & 0.368 & 46.7\,\%\\
\bottomrule
\end{tabular}%
}
\end{table}

\begin{figure}[!t]
\centering
\includegraphics[width=\columnwidth]{graphs/fig_us_training_method_bar.png}
\caption{IC by training protocol variant, S\&P~100. Moving from our
protocol~(A) to the standard Qlib protocol~(C) collapses IC by
84\,\%. Under Variant~C, hit rate falls to 46.7\,\%---below random---
indicating an inverted ranking signal.}
\label{fig:us_training_method_bar}
\end{figure}

Applying the Qlib standard (Variant~C) collapses IC from 0.043 to
0.007 (84\,\% reduction) and drops hit rate to 46.7\,\%, below the
50\,\% random baseline. Two mechanisms drive the failure:
\textit{DropExtremeLabel} removes tail observations that carry the
highest cross-sectional signal density, and removing BCE regularization
on the direction head destabilizes the shared embedding. Adding a
ranking term (Variant~D) does not recover the signal.

An important caveat: Variants~B--D were not HPO-tuned for their
respective loss functions. A fully fair comparison would require
separate search. However, a sixfold IC collapse and sub-random
directional accuracy are diagnostic of structural incompatibility, not
hyperparameter mismatch. MASTER and StockMixer use single-head designs
for which the Qlib protocol was developed. It does not transfer to
dual-head architectures. Researchers adapting benchmark training recipes
to architectures with different output structures should verify
compatibility before drawing conclusions.

\subsection{Cross-Market Analysis --- India (NIFTY~100)}
\label{sec:india}

No statistically significant cross-sectional signal is found on the
NIFTY~100. Table~\ref{tab:india_results} reports results for US
parameter transfer and India-specific HPO. A paired t-test yields
$p=0.574$ ($\Delta\mathrm{IC}=-0.001$); local optimization provides
no detectable improvement.

\begin{table}[!t]
\renewcommand{\arraystretch}{1.25}
\caption{India Market Results (NIFTY~100, $N=91$, Test Folds~4--7)}
\label{tab:india_results}
\centering
\resizebox{\columnwidth}{!}{%
\begin{tabular}{lcccc}
\toprule
\textbf{Configuration} & \textbf{IC (Mean$\pm$Std)} & \textbf{95\,\% CI}
                       & \textbf{Sharpe} & \textbf{Hit Rate}\\
\midrule
US HP Transfer & $0.017\pm0.033$ & $[-0.017,\;0.049]$ & 0.974 & 54.9\,\%\\
India HPO      & $0.016\pm0.026$ & $[-0.015,\;0.039]$ & 0.865 & 55.6\,\%\\
\midrule
\multicolumn{5}{l}{Paired t-test: $\Delta\mathrm{IC}=-0.001$,
$p=0.574$ (not significant)}\\
\bottomrule
\end{tabular}%
}
\end{table}

\begin{figure}[!t]
\centering
\includegraphics[width=\columnwidth]{graphs/fig_india100_ic_bar.png}
\caption{IC per fold, NIFTY~100 (India HPO). Fold~7 yields
IC~$=-0.031$; the mean IC of 0.016 produces
CI~$[-0.015,\,0.039]$ that spans zero.}
\label{fig:india100_ic_bar}
\end{figure}

\begin{figure}[!t]
\centering
\includegraphics[width=\columnwidth]{graphs/fig_india100_bootstrap_ic_hist.png}
\caption{Bootstrap IC distribution, NIFTY~100 (10\,000 resamples).
The lower CI bound is below zero.}
\label{fig:india100_bootstrap_ic_hist}
\end{figure}

\begin{figure}[!t]
\centering
\includegraphics[width=\columnwidth]{graphs/fig_us_vs_india_ic_ci.png}
\caption{Cross-market IC comparison with 95\,\% bootstrap CIs.
S\&P~100 CI lies entirely above zero; NIFTY~100 CI straddles zero
under both parameter configurations.}
\label{fig:us_vs_india_ic_ci}
\end{figure}

\begin{figure}[!t]
\centering
\includegraphics[width=\columnwidth]{graphs/fig_us_vs_india_per_fold_bar.png}
\caption{Per-fold IC, S\&P~100 (blue) vs.\ NIFTY~100 (red),
test folds~4--7. Both markets produce near-zero IC in Fold~6,
pointing to a shared global regime effect. US IC exceeds India IC
in every test fold.}
\label{fig:us_vs_india_per_fold_bar}
\end{figure}

Three structural conditions explain the null India result. First, with
$N=91$ and $k=20$, each stock connects to 22\,\% of the NIFTY~100
universe through the KNN module, approaching a fully connected graph
that eliminates the structural specificity the module is designed to
provide (the same ratio on the S\&P~100 is 20\,\%). Second, NIFTY~100
sector concentration is high: financial services accounts for roughly
35\,\% of index weight~\cite{nse2024nifty}, and several sectors contain
only one or two stocks, producing degenerate single-stock hyperedges in
the Static Sector HG---our ablation's most critical module. Third,
higher retail participation and T+1 settlement in Indian
markets~\cite{sebi2023settlement} introduce intraday noise patterns that
the 12 US-calibrated technical features do not capture. These conditions
define a \textit{market efficiency threshold}: the architecture requires
sufficient sector diversity, universe breadth, and microstructural
efficiency to generate propagatable cross-sectional signals.

\subsection{Benchmark Comparison}

Table~\ref{tab:benchmarks} places CRML against published architectures.
MASTER, StockMixer, and HIST were evaluated on Chinese A-share markets
(CSI~300/800) with 158-dimensional Alpha158 features; our primary
evaluation uses 12 technical indicators on the S\&P~100. A direct IC
comparison is informative but not valid given these differences. Our
primary empirical claim is architectural---that multi-scale relational
topology outperforms single-scale alternatives under identical
conditions.

\begin{table*}[!t]
\renewcommand{\arraystretch}{1.25}
\caption{Benchmark Contextualization. Sequential baseline IC values are
         from the Qlib benchmark leaderboard~\cite{yang2020qlib} for
         CSI~300 with Alpha158 features. Direct cross-market IC
         comparison is not valid due to differences in universe, market
         efficiency, features, and evaluation protocol.
         $^\dagger$Config~B-OUR: four-component loss, seq\_len$=60$;
         Config~B-8: MASTER protocol (MSE-only, CSZScoreNorm,
         DropExtreme, seq\_len$=8$).}
\label{tab:benchmarks}
\centering
\begin{tabular}{llccccc}
\toprule
\textbf{Model} & \textbf{Type} & \textbf{IC}
               & \textbf{Universe} & \textbf{Features}
               & \textbf{Year} & \textbf{Venue}\\
\midrule
\multicolumn{7}{l}{\textit{Sequential baselines}}\\
LSTM~\cite{hochreiter1997long}    & RNN       & $\sim$0.028 & CSI~300 & Alpha158 & 1997 & Neural Computation\\
GRU~\cite{cho2014learning}        & RNN       & $\sim$0.030 & CSI~300 & Alpha158 & 2014 & EMNLP\\
ALSTM~\cite{qin2017dual}          & RNN+Attn  & $\sim$0.036 & CSI~300 & Alpha158 & 2017 & IJCAI\\
SFM~\cite{zhang2017stock}         & Freq-RNN  & $\sim$0.032 & CSI~300 & Alpha158 & 2017 & KDD\\
\midrule
\multicolumn{7}{l}{\textit{Adaptation and routing models}}\\
TRA~\cite{lin2021learning}        & RNN+Route & $\sim$0.041 & CSI~300 & Alpha158 & 2021 & NeurIPS\\
DoubleAdapt~\cite{zhao2023doubleadapt} & Meta-IL & $\sim$0.044 & CSI~300 & Alpha158 & 2023 & KDD\\
\midrule
\multicolumn{7}{l}{\textit{Graph-based models}}\\
RSR~\cite{feng2019temporal}       & GNN       & --    & NYSE/NASDAQ  & Price      & 2019 & ACM TOIS\\
HATS~\cite{kim2019hats}           & Het-GNN   & --    & S\&P~500     & Price+Wiki & 2019 & arXiv\\
HIST~\cite{xu2021hist}            & Concept   & 0.048 & CSI~300/800  & Alpha158   & 2021 & arXiv\\
STHAN-SR~\cite{sawhney2021sthan}  & HGNN      & --    & NASDAQ       & Price      & 2021 & AAAI\\
MASTER~\cite{li2024master}        & Transf.   & 0.059 & CSI~300/800  & Alpha158   & 2024 & AAAI\\
StockMixer~\cite{fan2024stockmixer} & MLP-Mix & 0.055 & CSI~300/800  & Alpha158   & 2024 & AAAI\\
\midrule
\multicolumn{7}{l}{\textit{This work --- US S\&P~100 (primary)}}\\
\textbf{CRML (Proposed)} & \textbf{Multi-HG+GAT} & $\mathbf{0.047}$
  & \textbf{S\&P~100} & \textbf{12 Technical} & \textbf{2026} & \textbf{--}\\
Linear Baseline (Ridge) & Linear & $\sim$0.020 & S\&P~100 & 12 Technical & 2026 & --\\
\midrule
\multicolumn{7}{l}{\textit{This work --- CSI~300 cross-market
                   (Alpha158, not primary)$^\dagger$}}\\
CRML (Config~B-8) & Multi-HG+GAT & 0.016 & CSI~300 & 158 (Alpha158) & 2026 & --\\
CRML (Config~B-OUR) & Multi-HG+GAT & 0.041 & CSI~300 & 158 (Alpha158) & 2026 & --\\
\bottomrule
\end{tabular}
\end{table*}

CRML reaches IC~$=0.047$ on the S\&P~100 with 12~features, one unit of
IC below HIST's 0.048 on CSI~300/800 with 158~features. Given that the
S\&P~100 is one of the most efficiently priced equity universes in the
world and the feature set is $13\times$ smaller, this near-parity under
a strictly embargoed protocol is a useful data point.

The CSI~300 comparison isolates the protocol effect. Under the MASTER
protocol (Config~B-8), CRML reaches IC~$=0.016$---well below MASTER
(0.059) and StockMixer (0.055). Switching to our four-component loss and
longer sequence length (Config~B-OUR) recovers IC to 0.041, a $2.5\times$
improvement from protocol substitution alone. This directly corroborates
Section~\ref{sec:protocol}: the gap between CRML and CSI~300 benchmarks
is largely a protocol mismatch, not an architectural limitation.

\section{Conclusion}

This paper posed three research questions and the experimental results
provide direct answers.

\textbf{RQ1 is answered affirmatively.} Jointly modeling static sector
co-membership, dynamic similarity clustering, and directed pairwise
lead-lag generates a statistically significant cross-sectional signal on
the S\&P~100 with 12~technical features. Mean IC~$=0.047$ (bootstrap
CI~$[0.016,\,0.077]$) over four held-out test folds, against IC~$=0.020$
for ridge regression, confirms the gain is structural rather than
feature-driven.

\textbf{RQ2 is answered affirmatively.} Ablation shows that removing
any single module degrades performance. Sector hypergraph removal causes
the largest absolute IC drop ($\Delta\mathrm{IC}=-0.012$); KNN removal
yields the largest effect size (Cohen's $d=2.323$); GAT inclusion
raises long-short Sharpe by 13.5\,\% despite missing IC significance.
The gate equilibrium at 0.333 per module confirms complementary rather
than redundant contributions.

\textbf{RQ3 is answered negatively.} CRML does not transfer to the
NIFTY~100 under any tested configuration. Bootstrap CIs span zero for
both parameter transfer and India-specific HPO. Three structural
conditions---high KNN connectivity ratio relative to universe size,
degenerate single-stock sector hyperedges from index concentration, and
a US-calibrated feature set---define a market efficiency threshold below
which the architecture's relational propagation does not produce
detectable alpha.

A separate finding concerns training protocols. Applying the standard
Qlib MSE-only objective to CRML's dual-head design collapses IC by
84\,\% and drops directional accuracy below random. The BCE direction
loss is a structural component of the training objective. Benchmark
training recipes are not architecture-agnostic; compatibility should be
verified before reuse.

Four limitations bound these conclusions. First, four held-out test
folds limit statistical power; paired comparisons in the ablation do not
all reach conventional thresholds despite consistent IC improvements,
and Sharpe evidence partially fills this gap. Second, the 12-feature
input excludes fundamental data, analyst revisions, and alternative
signals that institutional models use. Third, the fusion gate converges
to uniform weights, suggesting the current gating design does not
realize the regime-dependent routing that Proposition~1 shows is within
the model's capacity. Fourth, the data covers one particular multi-year
period on two specific indices; the roll-forward window cannot
characterize all possible macro environments. Universe survivorship
bias~\cite{brown1992survivorship}---from using a 2026 index snapshot
rather than point-in-time membership---further qualifies generalizability.

Future work has several directions. Incorporating fundamental and
alternative data will test whether the relational topology amplifies
lower-frequency signals. Explicit symmetry-breaking regularization in
the gating layer may unlock the regime-dependent routing that the
architecture supports but does not consistently realize. Scaling to
broader universes such as the Russell~1000 will clarify the upper
boundary of the market efficiency threshold identified here, and
additional emerging market tests will clarify the lower boundary.

\bibliographystyle{IEEEtran}
\bibliography{references}

\end{document}